From - Mon May 22 18:39:58 2000
Received: from grumpy.usu.edu (grumpy.usu.edu [129.123.1.86])
	by math.usu.edu (8.9.3+Sun/8.9.1) with ESMTP id OAA07578
	for <symanzik@sunfs.math.usu.edu>; Tue, 16 May 2000 14:52:50 -0600 (MDT)
Received: from webmail.usu.edu ("port 3769"@webster.usu.edu [129.123.1.90])
 by cc.usu.edu (PMDF V5.2-32 #30472)
 with ESMTP id <01JPGZR6K8FU9LV17Q@cc.usu.edu> for symanzik@sunfs.math.usu.edu;
 Tue, 16 May 2000 14:54:36 MDT
Date: Tue, 16 May 2000 14:54:36 -0600
From: weipingdeng <weipingdeng@cc.usu.edu>
Subject: homework #5
Sender: weipingdeng <weipingdeng@cc.usu.edu>
To: symanzik@sunfs.math.usu.edu
Message-id: <39207441@webmail.usu.edu>
MIME-version: 1.0
X-Mailer: WebMail (Hydra) SMTP v3.61
Content-type: text/plain; charset="US-ASCII"
Content-transfer-encoding: 7bit
X-WebMail-UserID: weipingdeng
X-EXP32-SerialNo: 00002751
Content-Length: 2787
Status: RO
X-Mozilla-Status: 8001
X-Mozilla-Status2: 00000000
X-UIDL: 38baa60a000002e7

I explored the electronic textbook: Hyperstat Online at
http://www.davidmlane.com//hyperstat

This electronic textbook contains many contents. I think it is a kind of
combination of our several different classic textbooks in our different
courses. Chapter 4 (Introduction to Probability) is talked about in our Math
5710. Chapter 12, 13 and 14 are in our Stat 5200. And Chapter 15 (Prediction)
is in stat 5100. It is somewhat comprehensive.

I read Chapter 15 in detail. Most of the chapters have the same features as in
Chapter 15.

Chapter 15, just like our classic textbooks, was divided into several
sections. The contents are discussed step by step. It is well organized.After
we read this chapter, we can learn some concepts, statistics, models,
computations and so on. Only looking at the text inside Hyperstat, I felt this
is just ordinary text posted on the web, because there is no interactivity on
line that instructors and students can do by computer. But there are a lot of
links of references related to the chapter, including Analysis tools,
Instructional Demos, and Text. I'm interested in the Demos, so I clicked
Linear Regression. There is an applet that let us mark the locations of some
points (order pairs (X,Y)), and then the equation and graph of the regression
line will be given. If you change or add points, the regression line and
equation will change accordingly. This Demonstration is pretty good. It can
help students to understand what the regression line and equation are, and how
the points affect the line.

In class, Doctor Symanzik showed us a nice section in Chapter 5, Normal
Distribution. In section 5, click Z table and then we can get a three standard
normal Pdf graphs.. (a) If we type Z-value=-0.9 and click Compute Area, The
area below Z is given and the corresponding areas are painted blue in the
graphs immediately. (b) If we type 0.2 for Area below Z and click Compute Z,
then the Z value is given and the areas become blue. This part helps students
to visualize the relationship between P-value and Z-value.

This textbook includes Glossary, Search, Analysis tools, and Instructional
Demos. The contents are accurate, well organized. It has many examples,
references, pictures and case studies. It is a good textbook from traditional
view points.

In this book, each chapter has a lot of exercises, but students can't get
feedback immediately. Because this book on line covers so many different
issues, it's not so easy to read and understand. To understand each chapter,
we have to refer the related text. And as I said above, the interactivity on
line inside the Hyperstat is rare. I think this electronic textbook would be
better if it is improved further in some aspects according to the features of
computer.

Weiping Deng
05/15/00

